\begin{figure}[tbp]
  \centering
  \begin{tikzpicture}[
    x=1cm,y=1cm,
    letter/.style={minimum width=.48cm,minimum height=.48cm,inner sep=1pt},
    boundary/.style={letter,draw,line width=.6pt,rounded corners=1pt},
    block/.style={draw=black!55,fill=black!4,rounded corners=2pt,
      minimum width=3.25cm,minimum height=.68cm,inner sep=3pt},
    move/.style={-{Stealth[length=2mm]},line width=.65pt},
    note/.style={font=\small,align=center,inner sep=2pt}
  ]
    \node[note,anchor=west] at (-5.7,1.15)
      {A nonempty block and its boundaries};
    \node[boundary] (left) at (-4.7,0) {$e$};
    \node[block] (word) at (-2.45,0) {$p\;0\;\cdots\;p\;0$};
    \node[boundary] (right) at (-.2,0) {$e$};
    \node[letter] (end) at (1.0,0) {$w$};
    \draw (left) -- (word) -- (right) -- (end);
    \draw[decorate,decoration={brace,amplitude=4pt,mirror}]
      (-4.0,-.48) -- (-.9,-.48)
      node[note,midway,below=5pt] {letters in $\{p,0\}$};
    \node[note,align=left,anchor=west] at (1.75,0)
      {If the right $e$ is terminal,\\it may be removed.};

    \draw[move] (-2.45,-1.2) -- (-2.45,-1.85);
    \node[boundary] (leftnext) at (-4.7,-2.35) {$e$};
    \node[block] (wordnext) at (-2.45,-2.35) {$p\;0\;\cdots\;p\;0$};
    \node[boundary] (rightnext) at (-.2,-2.35) {$w$};
    \draw (leftnext) -- (wordnext) -- (rightnext);
    \node[note,align=left,anchor=west] at (1.75,-2.35)
      {The left $e$ is still not terminal:\\it cannot be removed first.};

    \draw[move] (-2.45,-2.95) -- (-2.45,-3.6)
      node[note,midway,right=4pt] {any further deletions};
    \node[note,anchor=west] at (-5.7,-4.0)
      {The final surviving block letter, with $q\in\{e,w\}$};

    \node[boundary] (ep) at (-4.7,-4.8) {$e$};
    \node[letter] (p) at (-3.5,-4.8) {$p$};
    \node[boundary] (qp) at (-2.3,-4.8) {$q$};
    \draw (ep) -- node[above,font=\small] {$1$} (p)
      -- node[above,font=\small] {$1$} (qp);
    \node[note] at (-3.5,-5.5) {$C$-weight: $1+1=0$};

    \node[boundary] (ez) at (.6,-4.8) {$e$};
    \node[letter] (zero) at (1.8,-4.8) {$0$};
    \node[boundary] (qz) at (3.0,-4.8) {$q$};
    \draw (ez) -- node[above,font=\small] {$0$} (zero)
      -- node[above,font=\small] {$0$} (qz);
    \node[note] at (1.8,-5.5) {$B$-weight: $0+0=0$};
  \end{tikzpicture}
  \caption{The obstruction in the second summand of the leading pentagon,
  after the minimal $A$-count forces every $e$ to be removed by an
  $A$-operation. Such an operation removes only a terminal $e$.
  While the displayed $p,0$ block is nonempty, its left boundary remains
  $e$; removing a terminal right boundary changes $e$ to $w$.
  With either right boundary, the final surviving $p$ or $0$ has deletion weight
  zero in $\mathbb F_2$. Thus no nonzero deletion sequence empties the block.}
  \label{fig:deletion-block}
\end{figure}
